\section{Introduction}

Binary reward --- assigning 1 if all test cases pass and 0 otherwise --- is the reward formulation used by all major published papers on reinforcement learning with execution feedback (RLEF) for code large language models (LLMs) \citep{le2022coderl,shojaee2023ppocoder,gehring2024rlef,yu2025dapo}. Yet under Group Relative Policy Optimization (GRPO), this widespread choice produces zero gradient contribution in the vast majority of early-training groups. In a 1,000-group simulation of realistic early-training conditions (partial-pass probability of 0.1 per test case, 8 completions per group, 5 test cases per problem), 987 of 1,000 groups (98.7\%) receive identical zero reward for every completion under binary reward --- zero group-mean, zero advantages, zero gradient. Ratio reward ($k/n$ test cases passing) escapes this dead zone in 98.7\% of those same groups, providing a mathematical guarantee of non-zero gradient signal whenever any completion achieves partial credit.

This gradient dead zone is not a probabilistic concern or a numerical artifact. It is a mathematical consequence of GRPO's group-normalization mechanics. GRPO computes advantages by subtracting the group mean from each completion's reward. When all completions in a group fail every test case --- the dominant scenario in early training on hard datasets --- binary reward assigns 0 to all completions, yielding a group mean of 0 and advantages of exactly 0 for every completion. No gradient flows. Ratio reward, by contrast, assigns $k/n$ for any completion that passes $k > 0$ test cases, creating non-zero advantage variance even in groups where no completion solves the problem completely. For a realistic 8-completion group with test-pass counts $[0, 0, 1, 2, 0, 3, 0, 0]$, binary reward gives all-zero advantages (variance: 0.0), while ratio reward gives advantages ranging from $-0.15$ to $+0.45$ (variance: 0.0475). The model learns to prefer the 3-test-pass completion over the 0-test-pass completions --- a learning direction that binary reward cannot provide.

Prior RLEF work has addressed reward sparsity at the policy level (curriculum learning, reward shaping) but has not analyzed the group-level dead zone that is specific to GRPO's normalization scheme. DAPO \citep{yu2025dapo} --- the most prominent GRPO-based code training work --- acknowledges reward sparsity concerns but uses binary reward and does not quantify the fraction of training groups that receive zero gradient contribution. CodeRL \citep{le2022coderl} and PPOCoder \citep{shojaee2023ppocoder} use PPO with a critic network that partially compensates for sparse rewards; GRPO eliminates the critic, making the dead zone more acute. No published work has conducted a controlled ablation of reward granularity under GRPO, nor has any work derived the exact condition under which binary reward must produce zero gradient contribution in a GRPO group.

We address this gap through a two-phase experimental design. In the first phase (h-e1), we prove the mechanistic claim: ratio reward provides a mathematical guarantee of non-zero advantage variance whenever the group contains at least one completion with $1 \leq k < n$ test cases passing and no completion passes all $n$ tests --- precisely the regime where binary reward guarantees zero variance. We confirm this via synthetic group validation, a 1,000-group early-training simulation, and a functional smoke test (3 GRPO steps, exit code 0). In the second phase (h-m1), we test whether this mechanistic difference translates to policy-level performance improvements in a controlled 1,000-step GRPO training run on APPS \citep{hendrycks2021apps} with DeepSeek-Coder-6.7B-instruct \citep{guo2024deepseek}. The h-m1 experiment produced a precision null result: both binary and ratio reward conditions showed \texttt{reward\_mean = 0.0} at all 208 logged training steps, with \texttt{clipped\_ratio = 1.0} throughout --- confirming that \texttt{max\_new\_tokens = 512} is insufficient for DeepSeek-Coder-6.7B to produce syntactically complete solutions for APPS competition problems, making both rewards mathematically equivalent (both = 0 for all completions). The null result is not a failed experiment: it precisely characterizes the experimental prerequisite for observing ratio reward's advantage --- the base model must achieve non-zero partial solve rates on the training dataset under the chosen generation length constraint.

We make the following contributions. \textbf{First}, we prove that ratio reward guarantees non-zero GRPO advantage variance in training groups where binary reward guarantees zero variance, and quantify that 98.7\% of early-training groups satisfy this condition under realistic partial-pass distributions (advantage variance: 0.0475 ratio vs.\ 0.0 binary; 987/1,000 simulation groups). \textbf{Second}, we provide the first quantitative analysis of binary reward's gradient dead zone in GRPO training for code LLMs, identifying the exact group-level condition under which the dead zone occurs. \textbf{Third}, we characterize the experimental prerequisite for observing policy-level ratio-vs-binary differences --- non-zero training-set solve rate --- and demonstrate that APPS competition-level problems with 512-token generation violate this prerequisite for a 6.7B-parameter model. \textbf{Fourth}, we validate a complete, reproducible GRPO + ratio reward training infrastructure (35/35 unit tests passing; 9 engineering issues documented and resolved), enabling future experiments under correctly-powered conditions.

The rest of this paper is organized as follows. Section~\ref{sec:related} discusses related work and its treatment of reward formulation. Section~\ref{sec:methodology} describes our experimental methodology. Section~\ref{sec:setup} details the experiments. Section~\ref{sec:results} presents results. Section~\ref{sec:discussion} discusses implications and limitations. Section~\ref{sec:conclusion} concludes.
